\section{Four ways to write the paper}\label{sec:designs}

Four designs were written out from the findings of Part~II. Each had to build only on what is solid or can be measured with data on disk, and to add no theory beyond repairs. Three reviewers scored them: one in the role of a theory editor at a top journal, one as an empirical editor, one as a senior co-author who must decide what the team does in the next three months. The ranking was unanimous.

\begin{center}\small
\setlength{\tabcolsep}{4.5pt}
\begin{tabularx}{\textwidth}{@{}L>{\centering\arraybackslash}p{0.45in}>{\centering\arraybackslash}p{0.7in}>{\centering\arraybackslash}p{0.62in}>{\centering\arraybackslash}p{0.72in}>{\centering\arraybackslash}p{0.9in}@{}}
\toprule
Design & One claim & Only solid ground & Evidence in hand & Feasible in three months & Reviewers' odds at a top-5 journal \\
\midrule
\textbf{1. Scale versus speed.} More researchers widen science; faster researchers thicken it first. & 8 & 8 & 5 & 7 & 4--10\% \\
2. The same, plus ``crowding costs about 1\%'' as a second claim. & 6--7 & 6 & 5 & 7 & 3--8\% \\
3. Incidence. Who gains from research tools? & 6 & 6--7 & 3 & 5--6 & 2--6\% \\
4. \emph{Too Fast to Read} as a dynamic paper on catch-up and overshooting. & 5--6 & 4--5 & 1--2 & 3--4 & 1--2\% \\
\bottomrule
\end{tabularx}
\par\smallskip
\parbox{\textwidth}{\footnotesize Scores out of 10, range across the three reviewers. The odds are for the best case. They are judgements, not estimates, and they are the reviewers' own: where the designer of a paper had given a higher figure, the reviewers lowered it.}
\end{center}

\begin{center}\small
\begin{tabularx}{\textwidth}{@{}P{1.0in}LL@{}}
\toprule
Design & Strongest point & Weakest point \\
\midrule
1.\ Scale versus speed & One contrast a reader keeps, built on the two best facts. Its new half can fail in two directions with data on disk. It uses nothing the review refuted, and if the test fails the authors have lost about five weeks. & The robust half of the test (extra papers stay in their authors' topics) is predicted by any topic-specific skill. The discriminating half (entry falls) is not pinned for 2026 and may not be visible until 2027. \\
2.\ With the 1\% claim & The best-ordered plan (two gates) and the clearest table of what each view predicts. & Two claims. The 1\% rests on Lemma~4, on a slope from a regression that cannot be reproduced, and on no fertility depletion. \\
3.\ Incidence & The best economic idea: a gain that reaches only entrants is capitalised into the entry barrier. An honest repair of the central theorem. & The headline cannot fail and cannot be tested before 2027. After the repair it is a band, from a third to nine tenths of two log points, set by inputs that are unmeasured or assumed. \\
4.\ Dynamic & The most original mechanism: the entry cost is a stock, and stocks lag flows. Every object is observable. & No evidence today. Its lead result rests on a law of motion without the catch-up delay. It keeps a title whose object it has to demote. \\
\bottomrule
\end{tabularx}
\end{center}

The five readings of the first round had put incidence first. The three reviewers disagreed, and the reason deserves to be stated plainly. The readings ranked what was in the manuscript, and the 2026 shock was not in it. Incidence is a framework: its statistics are inputs, and nothing in the data can contradict it this year. The contrast between scale and speed is a claim, and the data on disk can reject it.

\section{The recommended paper}\label{sec:recommended}

\subsection{Title, question, thesis}

\begin{keybox}[The paper]
\textbf{Working title.} \emph{More Researchers or Faster Researchers? How Science Absorbs Growth.} To be adopted only if the test supports it.

\textbf{Question.} Does science absorb a rise in output per researcher the way it absorbs a rise in the number of researchers?

\textbf{Thesis.} Newcomers are the margin on which science adjusts, because joining a topic has an entry cost that incumbents have already paid. So more researchers open new topics, while faster researchers first thicken the topics they are already in.
\end{keybox}

\needspace{15\baselineskip}
\subsection{The abstract it could have}

Brackets mark what is not yet established.

\begin{quote}\small
Science can grow through more researchers or through faster ones. The two are absorbed differently. Researchers choose topics, and joining a topic means catching up on its literature. Incumbents have paid that cost and earn a rent, so they stay; newcomers are the margin that adjusts. More researchers therefore open new topics, while faster researchers thicken the topics they are in until entry responds. On arXiv, growth in 2018--2025 came from new authors, and local density rose with an elasticity of [0.12--0.19], against [0.53--0.81] in a fixed idea space. In January--September 2026 papers per author rose 11--14\% after seven flat years. [The extra papers landed in their authors' existing topics, which thickened; new authors again opened new ones.] The model predicts that entry into those topics falls as their literature builds up; we state this as a dated prediction. The welfare effects are small.
\end{quote}

\subsection{The results it would contain}

\begin{center}\small
\begin{tabularx}{\textwidth}{@{}P{0.95in}LP{1.35in}@{}}
\toprule
 & Statement & Comes from \\
\midrule
\textbf{Scale} & More researchers open new topics at once; local density moves with their number at a small elasticity. & Proposition 3, Remark 4, Fact 1 (repaired) \\
\textbf{Speed} & A rise in output per researcher raises the flow of every existing topic. Insiders stay, and the excess is worked off as fast as members turn over. In the long run speed and scale are the same. & Appendix B.8 and Theorem 2(iv), rewritten as one proposition; the 2026 surge \\
\textbf{Why they differ} & Insiders earn a premium equal to the newcomer's entry cost, and their payoff rises with density where the newcomer's falls. & Proposition 5, Corollary 1, Lemma 2 \\
\textbf{The comparison} & If nobody had to pay to get in, speed and scale would be absorbed alike. The asymmetry is the footprint of the entry cost. & A five-line proposition for the frictionless case \\
\textbf{Tools} & Inside a topic a tool is worth the time it saves. A tool that makes researchers faster crowds existing topics first. A tool that makes entry cheaper raises density and insiders' returns. & Lemma 5; the sign statements of \S\ref{sec:ideas-incidence} \\
\bottomrule
\end{tabularx}
\end{center}

The rents are the \emph{mechanism} of the contrast, not a third result. The tools are an \emph{implication}, stated as signs with ranges and no ranking. If either grows back into a headline the paper is several papers again.

\subsection{Outline}

\begin{center}\small
\begin{tabularx}{\textwidth}{@{}P{0.25in}P{1.9in}LP{0.5in}@{}}
\toprule
 & Section & Content & Pages \\
\midrule
1 & Introduction & The question; the two facts side by side; the mechanism in one paragraph; what is and is not claimed. & 3 \\
2 & Two facts & Fact A: the repaired volume elasticity, with its benchmark at every radius. Fact B: the 2026 rise in papers per author. & 5 \\
3 & Model & Topics, the payoff curve (stated; proof in the appendix), the exact catch-up cost, the frontier value, equilibrium. Fewer than 20 symbols. & 5 \\
4 & Scale and speed in the model & The five results above. & 5 \\
5 & The 2026 test & Measurement; where extra papers land; crowding and entry; threats. & 8 \\
6 & What this means for research tools & Lemma 5 and the two sign statements. & 2 \\
7 & Magnitudes & The two elasticities, turnover, adjustment time; the small numbers as findings. & 2 \\
8 & Conclusion & The dated predictions for 2027. & 1 \\
\midrule
 & Appendices & Proofs; data and name keys; robustness; the rule changes as context; the calibration as an illustration. & \\
\bottomrule
\end{tabularx}
\end{center}

About thirty pages. The reviewers' advice was to aim lower still.

\subsection{What goes where}

\begin{center}\small
\begin{tabularx}{\textwidth}{@{}LP{2.2in}@{}}
\toprule
Material in the current draft & Place \\
\midrule
Proposition 3, Remark 4, Fact 1 & Lead: scale \\
Appendix B.8, Theorem 2(iv), the 2026 surge & Lead: speed \\
Proposition 5, Corollary 1 (signs), Lemma 2 without the label ``speed limit'' & The mechanism \\
Lemma 5; the repaired output accounting & A short section on tools \\
Theorem 3 with the inframarginal term; the founders' share & The fallback, or a second paper \\
Theorem 1; the calibration; Propositions 4 and 10 as ranges & Appendix, plus a paragraph of magnitudes \\
Theorem 2(i)--(iii); Propositions 9 and 11; the life-cycle simulation & Shelved, for a paper with its own evidence \\
Lemma 1; Facts 3--5; the listing designs; the cap; Propositions 6 and 8 & Out, or a data appendix \\
Fact 2 & Kept only as an upper bound on the catch-up gradient \\
Fifteen predictions & Three or four, about the path in 2027 \\
``What we no longer claim''; every sentence about earlier versions & Deleted \\
\bottomrule
\end{tabularx}
\end{center}

\section{The tests}\label{sec:tests}

\subsection{Two facts}

\textbf{Fact A (scale).} Repair Fact~1 on its own panel: standard errors clustered by area or by wild bootstrap, every radius next to its benchmark, the years before 2023, embedding-based neighbours. Then build a version that runs from the repository: arXiv's primary categories, density measured as a flow (neighbours posted in the previous twelve months), regressed on the number of authors and on papers per author separately. Both elasticities of the paper must come from one panel and one measure.

\textbf{Fact B (speed).} Write up the 2026 rise by field and by month, with authors and not submitters, for a fixed cohort, with papers per name capped at three, five and ten, with and without catch-all categories. This needs no text measure and can be done first.

\subsection{Three tests}

\begin{center}\small
\begin{tabularx}{\textwidth}{@{}P{1.05in}LL@{}}
\toprule
Test & Construction & What supports, what refutes \\
\midrule
\textbf{A. Where do the extra papers land?} & A fixed cohort of author names active in each of 2022--2024. For each of their papers: is it in a topic cell where the author published before? Is its category one of the author's earlier categories (a version that uses no text)? How many papers by others lie nearby? Regress on the author's change in output, with author effects and category-by-year effects. Placebo: the change from 2024 to 2025. & \emph{Supports:} authors who sped up keep the same share of papers in their home cells, and the extra papers sit in cells at least as dense. \emph{Refutes:} the extra papers land outside home cells, in thin or new ones. Then speed behaves like scale. \\
\textbf{B. Do their topics thicken, and does entry respond?} & Topic cells fixed on 2022--24 text, at three sizes chosen in advance. Split each cell's flow into papers by incumbents, by movers (already on arXiv, new to the cell) and by first-time authors. The cell's speed shock is the change in its incumbents' output \emph{outside} the cell, so that an opportunity specific to the cell does not drive it. Regress the change in the cell's flow on the speed shock and on growth in authors; regress movers' entry on the speed shock. & \emph{Supports:} the flow loads on the speed shock clearly more than on growth in authors. \emph{Half a result:} flow up, entry unchanged; the fall in entry is then the dated prediction. \emph{Refutes:} speed and scale load alike. \\
\textbf{C. Two elasticities in one regression} & Categories by half-year, 2019--2026: local density on the number of authors and on papers per author. & A consistency check only. Across categories the rise in papers per author has a standard deviation of 0.04, about 0.03 of it signal, so the standard error is near 0.3. \\
\bottomrule
\end{tabularx}
\end{center}

The cell-level and author-level tests carry the paper. The cohort is large (of the order of 100,000 names), which is where the power is.

\subsection{Threats, and what answers them}

\begin{center}\small
\begin{tabularx}{\textwidth}{@{}P{1.45in}L@{}}
\toprule
Threat & Defence \\
\midrule
\textbf{The wording of abstracts changed in 2025--26.} Word-based similarity fell; embedding-based similarity may err the other way, and the authors who sped up may be the ones whose wording changed. & Comparisons within a category and period. Versions that use no text: the author's home category; cells built from shared references in high-energy physics. A ``style floor'': the similarity of random pairs across categories in each half-year, with radii defined relative to it. A subsample of abstracts rewritten by one fixed procedure before embedding. \\
\textbf{Three rule changes in the same months.} & Count authors, never submitters. Define displaced entrants as authors already on arXiv, on whom the endorsement rule does not bind. Drop review-titled papers. Make mathematics and physics the main sample. End the window on 30 September, before the cap. \\
\textbf{The extra papers may be of low quality.} If newcomers ignore them, they add nothing to what must be read. & A fixed cohort of established authors. Drop catch-all categories and withdrawn papers. In high-energy physics, split the extra papers by early citations. This is the hardest threat: early citations are thin. \\
\textbf{The rise is concentrated among prolific names.} & Caps at three, five and ten papers per name; rare surnames; names above the 99th percentile trimmed. \\
\textbf{A common shock raises output and density together.} & The leave-out construction of the speed shock. Flat pre-trends in 2024--25. A lead test: the size of a cell's 2026 shock must not predict its change in 2024--25. \\
\textbf{The result is read after the fact.} & A dated memo, deposited before any 2026 similarity is computed: outcomes, samples, radii fixed on 2019--22, cells fixed on 2022--24, thresholds for ``supports'' and ``refutes'', and a power calculation on the placebo years. \\
\bottomrule
\end{tabularx}
\end{center}

\subsection{Is the entry cost the literature?}

None of the three tests shows that the cost of entry is \emph{reading}. Two by-products of the same panel bear on it. First, measured turnover: the share of a cell's papers written by authors new to it, by year, which replaces the assumed rate of exit and gives the adjustment time. Second, the barrier signature of \S\ref{sec:ideas-other}: holding a cell's recent flow fixed, does the newcomers' share fall with the stock of older papers, while incumbents' returns rise and their exit does not? If the stock gradient is absent, the paper says ``a sunk entry cost'' and keeps Lemma~2 as its leading example.

\needspace{12.5cm}
\subsection{Two gates}

\begin{figure}[h]
\centering
\begin{tikzpicture}[font=\sffamily\small,
 gate/.style={draw=hblue,line width=0.8pt,rounded corners=2pt,fill=hblue!6,align=center,inner sep=6pt,text width=4.3cm},
 outc/.style={draw=axisc,line width=0.6pt,rounded corners=2pt,fill=white,align=left,inner sep=6pt,text width=6.9cm},
 stopc/.style={outc,draw=stcrit,fill=stcrit!5},
 arr/.style={-{Stealth[length=5pt]},line width=0.8pt,color=inktwo},
 lab/.style={font=\sffamily\footnotesize,color=inktwo}]
\node[gate] (g1) at (0,0) {\textbf{Gate 1}\\Does the main fact survive on a panel a referee can re-run?};
\node[gate] (pre) at (0,-2.7) {\textbf{Before looking}\\Derive the model's prediction for 2026 and deposit it};
\node[gate] (g2) at (0,-5.5) {\textbf{Gate 2}\\Where did the extra papers of 2026 land?};
\node[stopc] (s1) at (7.6,0) {\textbf{No.} No version of the paper stands. What remains is Lemma~5, the rent results and the 2026 fact, as a short note.};
\node[outc] (o1) at (7.6,-3.35) {\textbf{In their authors' topics, which thickened.} Write \emph{scale versus speed}. State the fall in entry as a dated prediction, and submit after a first reading of 2027.};
\node[outc] (o2) at (7.6,-5.5) {\textbf{Elsewhere: speed looks like scale.} Write the fallback: the entry margin, rents, and incidence with a measured founders' share.};
\node[outc] (o3) at (7.6,-7.65) {\textbf{The measures disagree.} Report the surge as a fact with the registered prediction, write the fallback, and return with 2027 data.};
\draw[arr] (g1) -- node[lab,right] {yes} (pre);
\draw[arr] (pre) -- (g2);
\draw[arr] (g1) -- (s1);
\draw[arr] (g2.east) -- ++(0.55,0) |- (o1.west);
\draw[arr] (g2.east) -- (o2.west);
\draw[arr] (g2.east) -- ++(0.55,0) |- (o3.west);
\end{tikzpicture}
\caption{The order of decisions. The first month of work is the same under every outcome, so nothing is lost by choosing the headline at the end of it.}
\label{fig:gates}
\end{figure}

\subsection{The fallback and the second paper}

\textbf{The fallback.} If speed looks like scale, incumbents are mobile at the margin, and every result that needs members to stay put is rejected with it (Appendix~B.8, Theorem~2, Proposition~9). The paper is then the one the first round agreed on: the entry margin (Proposition~3 and Fact~1, with a second fact that a rise in output per author did not thicken topics either), rents as the reading of the returns gradient, and the incidence of tools with sufficient statistics. That paper needs the founders' share.

\textbf{Measuring the founders' share.} In high-energy physics, where reference lists are complete: call a paper a \emph{pioneer} if no earlier paper lies within a fixed radius in the previous 24 months, and its cell \emph{founded} if at least a given number of later papers by other authors land in it. For the authors of pioneers, split their citation-weighted output over the next five years into work in cells they founded and that drew followers, work in cells founded by others, and work that nobody followed. The founders' share is roughly the first as a fraction of the first and the third. Two to three weeks. The radius is the risk: fix it and the thresholds in the dated memo, and report across radii. The same panel gives the share of a followed cell's output written by early joiners.

\textbf{The second paper.} Incidence becomes a strong paper once the entry mechanism has been shown in data and a shock on the joiners' side has outcomes. The candidate is the review rule in computer science, with exposure measured from what newcomers cited before the rule, and outcomes through 2027.

\textbf{The dynamic paper.} Do not start it. Revisit it only if the stock regressions are strong and data on the review rule exist.

\subsection{What 2027 must show}

If the contrast is found in 2026, the paper should state in advance what the next year of data must look like:
the excess density of the cells that thickened is closing;
entry by movers into those cells is below replacement;
and the new authors of 2026 again opened new cells, at the elasticity of Fact~1. The last is a check on Fact~1 itself: if the new authors of 2026 thickened existing cells, the main fact fails out of sample, and with it every version of the paper.
